\section{Conclusion}

Nghiên cứu này cho thấy động lực học PPG trên các cửa sổ ngắn thay đổi có hệ thống trong quá trình chuyển từ Awake sang Drowsy. Pseudoperiodic surrogate testing cho thấy ở cả hai trạng thái, tổ chức theo thời gian của tín hiệu gốc không được tái tạo đầy đủ bởi một noisy pseudoperiodic process, nhưng kết quả này không được diễn giải như bằng chứng của deterministic chaos. Trong phân tích chính với cửa sổ 60~s, Drowsiness đi kèm với giảm short-term forecastability, giảm tổ chức recurrence theo đường chéo và giảm local trajectory divergence, được phản ánh lần lượt bởi $CC\downarrow$, $NRMSE\uparrow$, $DET\downarrow$ và $LLE\downarrow$. Các thay đổi chính này duy trì cùng direction trên các cửa sổ 30--180~s.

Nhìn tổng thể, Awake--Drowsy transition không được đặc trưng bởi một thay đổi đơn chiều của ``chaos'' hay ``complexity'', mà bởi sự tái tổ chức phối hợp của nhiều thuộc tính động lực học bổ sung của PPG. Trong bối cảnh wake-to-sleep transition, pattern này phù hợp với sự thay đổi của peripheral cardiovascular dynamics trong một physiological regime đang chuyển tiếp và chưa ổn định. Mặc dù nghiên cứu chưa xây dựng classifier và chưa xác lập causal physiological mechanisms, các kết quả cho thấy short-window PPG chứa thông tin động lực học có thể được đặc trưng một cách nhất quán và tạo cơ sở cho các nghiên cứu tiếp theo về nonlinear physiological descriptors, multimodal validation và wearable drowsiness monitoring.